\documentclass[11pt]{article}
\usepackage[margin=1in]{geometry}
\usepackage{amsmath,amssymb,amsthm,mathtools}
\usepackage{graphicx}
\usepackage{microtype}
\usepackage[colorlinks=true,linkcolor=blue,urlcolor=blue]{hyperref}

\newtheorem{theorem}{Theorem}
\newtheorem{lemma}{Lemma}
\newcommand{\D}{\mathbb D}
\newcommand{\T}{\mathbb T}
\newcommand{\C}{\mathbb C}
\newcommand{\Hh}{P^2(\mu)}
\renewcommand{\thesection}{\arabic{section}.}

\title{Polynomial and Hardy-equivalent cyclicity in $P^2(\mu)$}
\author{A candidate substantial partial answer to arXiv:1809.10182}
\date{13 August 2026}

\begin{document}
\maketitle

\noindent\fbox{\parbox{0.94\linewidth}{\textbf{Status: candidate substantial
partial result, likely valid.}  Under the full source hypotheses, every
finite-codimensional shift-invariant subspace is cyclic, and these are exactly
the invariant subspaces containing a nonzero polynomial.  The full question is
also affirmative for Hardy-equivalent measures, including every rotationally
invariant example.}}

\section{Source question and result}

Let $\mu$ be a finite positive measure supported on $\overline\D$, let
$\Hh$ be the closure of the analytic polynomials in $L^2(\mu)$, and let
$S f=zf$.  Question~2.1 of Akeroyd--Conway--Yang [1] assumes that $\Hh$ is
irreducible, $\operatorname{abpe}(\Hh)=\D$, and $\mu(\T)>0$, and asks whether
every nonzero closed $S$-invariant subspace is cyclic.

\begin{figure}[h]
 \centering
 \includegraphics[width=0.88\linewidth]{figures/source_question_crop.png}
 \caption{The exact cyclicity question on page 3 of arXiv:1809.10182.}
\end{figure}

For a polynomial $q$, let $q_{\D}$ be the monic polynomial obtained by retaining
with multiplicity precisely the zeros of $q$ in $\D$; it is $1$ when there are
none.  Write $[f]=\overline{\{pf:p\in\C[z]\}}^{\,\Hh}$.

\small

\begin{theorem}[finite-codimensional classification]\label{thm:class}
Under the hypotheses above, for every nonzero polynomial $q$,
\[
 [q]=q_{\D}\Hh,
 \qquad \operatorname{codim}[q]=\deg q_{\D}.
\]
For a closed $S$-invariant subspace $M\ne\{0\}$, the following are equivalent:
\begin{enumerate}
 \item $M$ contains a nonzero polynomial;
 \item $M$ has finite codimension in $\Hh$;
 \item $M=d\Hh$ for a monic polynomial $d$ all of whose zeros lie in $\D$.
\end{enumerate}
In particular every such $M$ is cyclic, with generator $d$.
\end{theorem}

\begin{theorem}[Hardy-equivalent measures]\label{thm:hardy}
Write $\mu|_{\T}=h\,m$ and $\nu=\mu|_{\D}$.  If
\[
 0<c\le h\le C<\infty\quad(m\text{-a.e.})
\]
and $\nu$ is an $H^2$-Carleson measure, then every closed invariant subspace
of $\Hh$ is cyclic.  In particular this holds for every rotationally invariant
$\mu$ satisfying the source hypotheses.
\end{theorem}

\section{Boundary evaluations and division}

The first point is that irreducibility removes a possible boundary obstruction.

\begin{lemma}[no boundary bounded point evaluations]\label{lem:boundary}
$\operatorname{bpe}(\Hh)\cap\T=\varnothing$.  Consequently $S-\zeta$ has
dense range for every $\zeta\in\T$.
\end{lemma}

\begin{proof}
Fix $\zeta\in\T$ and set
\[
 u_n(z)=\left(\frac{1+\overline\zeta z}{2}\right)^n.
\]
On $\overline\D$, $|u_n|\le1$ and $u_n$ tends pointwise to
$\mathbf1_{\{\zeta\}}$.  If $\mu(\{\zeta\})=0$, dominated convergence gives
$\|u_n\|_{\Hh}\to0$, whereas $u_n(\zeta)=1$, ruling out bounded evaluation.
If $\mu(\{\zeta\})>0$, then $u_n\to\mathbf1_{\{\zeta\}}$ in $L^2(\mu)$.
This puts a nontrivial characteristic function in $\Hh$, contrary to
irreducibility.

If $S-\zeta$ did not have dense range, some nonzero $v$ would satisfy
$S^*v=\overline\zeta v$.  Then
$\langle p,v\rangle=p(\zeta)\langle1,v\rangle$ for polynomials $p$.
Here $\langle1,v\rangle\ne0$, since otherwise density of the polynomials would
give $v=0$.  Thus evaluation at $\zeta$ would be bounded, a contradiction.
\end{proof}

\begin{lemma}[bounded interior division]\label{lem:division}
For $a\in\D$, the map
\[
 D_a p=\frac{p-p(a)}{z-a}
\]
extends boundedly to $\Hh$.  Hence
\[
 (S-a)\Hh=\ker(\operatorname{ev}_a).
\]
More generally, if all zeros of $r\in\C[z]$ lie in $\D$, then $r\Hh$ is
closed, has codimension $\deg r$, and is described by vanishing of the
corresponding zero jets.
\end{lemma}

\begin{proof}
Choose $0<\rho<R$ with $\overline{B(a,R)}\subset\D$.  Analytic bounded point
evaluations are locally uniformly bounded, so
\[
 K=\sup_{|w-a|=R}\|\operatorname{ev}_w\|<\infty.
\]
The maximum principle applied to $D_ap$ gives, on $B(a,\rho)$,
\[
 |D_ap|\le R^{-1}(K+\|\operatorname{ev}_a\|)\|p\|_{\Hh}.
\]
On its complement, $|z-a|\ge\rho$ and therefore
\[
 |D_ap(z)|\le\rho^{-1}(|p(z)|+|p(a)|).
\]
Integrating these two estimates proves
$\|D_ap\|_{\Hh}\le C_a\|p\|_{\Hh}$.  Completion and locally uniform
convergence identify the extension with analytic divided difference.  Thus
every $f$ with $f(a)=0$ equals $(z-a)D_af$, proving the displayed range
identity.  Iteration, boundedness of derivative evaluations, and Hermite
interpolation give the last assertion and the exact codimension.
\end{proof}

\section{Proof of the classification}

\begin{proof}[Proof of Theorem~\ref{thm:class}]
Factor $q=q_{\D}q_{\T}q_{\rm out}$ according as its roots lie in $\D$, on
$\T$, or outside $\overline\D$.  Since $\|S\|\le1$, every exterior factor is
invertible.  Lemma~\ref{lem:boundary} says each boundary factor has dense
range; the commuting product of these factors also has dense range (equivalently,
the product of their injective adjoints is injective).  Since
$q_{\D}\Hh$ is closed by Lemma~\ref{lem:division}, continuity gives
\[
 [q]=q_{\D}\Hh.
\]
The codimension assertion follows from the same lemma.

Now suppose $q\in M\cap\C[z]$ is nonzero.  Then
$q_{\D}\Hh=[q]\subset M$, so $M$ has finite codimension.  Modulo
$q_{\D}\Hh$, polynomials exhaust the finite-dimensional cyclic module
\[
 \Hh/q_{\D}\Hh\cong\C[z]/(q_{\D}).
\]
The image of $M$ is a submodule.  Its inverse image in $\C[z]$ is an ideal
containing $(q_{\D})$, hence equals $(d)$ for a monic divisor $d$ of
$q_{\D}$.  Because $d\in M$, invariance and closedness give
$d\Hh\subset M$; the quotient description gives the reverse inclusion.
Thus $M=d\Hh$.

Conversely, if $M$ has finite codimension, let $Q$ be the operator induced by
$S$ on $\Hh/M$.  Its minimal polynomial $q$ satisfies $q\Hh\subset M$, so
$q\in M$.  The preceding argument applies.  (The eigenvalues of $Q^*$ are
bounded point evaluations; contractivity places them in $\overline\D$ and
Lemma~\ref{lem:boundary} excludes $\T$.)  Finally, $d\Hh=[d]$ is cyclic.
\end{proof}

\section{The Hardy-equivalent class}

\begin{proof}[Proof of Theorem~\ref{thm:hardy}]
If $K_\nu$ is the Carleson embedding constant, then for every polynomial $p$,
\[
 c\|p\|_{H^2}^2
 \le \|p\|_{L^2(\mu)}^2
 \le (C+K_\nu)\|p\|_{H^2}^2.
\]
Thus $\Hh$ and $H^2$ are the same topological shift module.  Their closed
invariant subspaces coincide, and Beurling's theorem writes each nonzero one
as $\theta H^2=[\theta]$.

If $\mu$ is rotationally invariant, its nonzero boundary part is $c,m$.
For the radial interior part and $p(z)=\sum a_nz^n$, angular orthogonality gives
\[
 \int_{\D}|p|^2\,d\nu
 =\sum_n|a_n|^2\int_{\D}|z|^{2n}\,d\nu
 \le\nu(\D)\|p\|_{H^2}^2.
\]
So the hypotheses just proved are automatic.
\end{proof}

\section{Scope, novelty, and review focus}
\small
Theorem~\ref{thm:class} removes all algebraic and finite-codimensional cases,
and Theorem~\ref{thm:hardy} handles a natural infinite-dimensional class.  The
remaining regime is non-Hardy-equivalent, infinite-codimensional, and
polynomial-free; index one alone does not bridge this gap.

A bounded exact-title and keyword search on 13 August 2026 found no later
solution or this classification.  Novelty confidence is modest because the
proof uses short module-theoretic arguments.

\medskip\noindent\textbf{Human-review focus.}
Check bounded divided differences in Lemma~\ref{lem:division}, the passage from
dense boundary factors to $[q]=q_{\D}\Hh$, and the finite cyclic-module ideal
argument.  Do not promote this as a full answer without treating
polynomial-free invariant subspaces.

\medskip
\noindent{\large\bfseries References}\par\smallskip
\footnotesize
\noindent[1] J.~R. Akeroyd, J.~B. Conway, and L.~Yang,
\emph{On Nontangential Limits and Shift Invariant Subspaces},
arXiv:1809.10182 (2018).\par
\noindent[2] A.~Aleman, S.~Richter, and C.~Sundberg,
\emph{Nontangential limits in $P^t(\mu)$-spaces and the index of invariant
subspaces}, Ann. of Math. 169 (2009), 449--490.

\end{document}
